%! Unit 3, Lessons 3.1–3.5 — ACTUAL Unit Test (Answer Key)
% The scored unit test. Item for item it is the parallel form of ../practice_test/: same two
% sections, the same ten multiple-choice targets in the same order, the same two free-response
% questions with the same parts and point values — every context different, so a student who
% worked the practice test has practised the moves and not memorised the answers.
% BLUEPRINT (1 pt per MC, FRQ 1 = 4, FRQ 2 = 6, total 20; one 55-minute period):
%  MC1  3.1 observational vs experiment — careful measurement is not a treatment (the crux)
%  MC2  3.1 generalize only to the population sampled
%  MC3  3.2 name the sampling method (cluster here; stratified on the practice test)
%  MC4  3.2 equal chance for each individual is not an SRS (one whole row at random)
%  MC5  3.3 a random sample can still be biased (response bias: leading wording)
%  MC6  3.3 a larger sample does not fix bias (voluntary response)
%  MC7  3.4 why random assignment: balance confounding variables
%  MC8  3.4 placebo and blinding (single-blind here; double-blind on the practice test)
%  MC9  3.4 matched pairs — chance still decides the order for every subject
%  MC10 3.5 scope of inference: random selection, no random assignment
%  FRQ1 3.2/3.3 SRS by random number generator; cluster vs stratified; nonresponse + direction
%  FRQ2 3.4/3.5 units/treatments/response; random assignment; blocking; reading a reshuffle
%       dotplot (NOT significant here, unlike the practice test); cause and to whom.
% THE DATA
%  FRQ2(d) 100 reshuffles, difference in mean yield (kg, new - standard), step 0.25:
%    -1.25:1 -1.0:1 -0.75:5 -0.5:11 -0.25:20 0:24 0.25:20 0.5:11 0.75:5 1.0:1 1.25:1
%    (sum 100). Observed 0.5: 11+5+1+1 = 18 of 100 at or above -> not significant.
% NAMESTRIP: no name line — mirrors the blank (a graded page; the name is on the cover sheet only).
% NO teachernote here — scoring guidance lives in ../../rubrics/, never in a key.
% GENERATED FROM the blank: the correct MC options wrapped in \ans{}, each
% \gradedhide{...\writelines{n}} turned into the box plus n \ansline{}.
% PAGINATION: each item carries a \boxguard sized to its height so no stem or option list
% splits across a page; the Section II strip is guarded to travel with FRQ 1.
% NO SKETCH ITEMS: the one display (the reshuffle dotplot) is pre-drawn in TikZ and read.
% PAGE LOCKSTEP with ../../test_keys/actual_test_key/: every \writelines{n} here is answered by
% exactly n \ansline{} there, each terse enough not to wrap. Nothing else differs except the
% preamble, the header, and the \ans{} marks on the correct MC options.
\documentclass[12pt]{article}
\usepackage{apstats-article}
\usepackage{apstats-key}   % pulls in -boxes; do NOT also load -boxes
\usepackage{grading}
\gradingassignment{apstats-u03-actual-test}
\gradingcourse{advanced_statistics}

\begin{document}

\pageheader{Unit 3 \quad Lessons 3.1--3.5}{Unit Test --- Answer Key}
{\small\itshape Write each final answer in its box. Only what is inside a box is graded.}\par

\hfill\small Score:\;\blank{1.3cm}\;/\;20\normalsize
\vspace{0.06in}

\boxguard[10]
\begin{headlinebox}{goldbg}
\small
\textbf{Section I --- Multiple Choice} (1 point each). Write the letter of the single best answer in its Letter box. Read all
five options --- several name a real idea from the unit, but only one fits \emph{this} study.
\end{headlinebox}

\vspace{0.10in}

\begin{enumerate}[leftmargin=1.9em, itemsep=0.13in]

  \boxguard[11]
  \item A nutrition team gave 250 office workers a food-logging app that recorded every meal they
        ate for three months, and weighed each worker every week. Workers who ate more fast food
        tended to gain more weight. Which statement is correct?
        \begin{enumerate}[label=\textbf{(\Alph*)}, leftmargin=2.2em, itemsep=2pt]
          \item It is an experiment, because every meal was recorded precisely.
          \item \ans{It is an observational study, because the team did not assign what anyone ate.}
          \item It is an experiment, because the team gave every worker the app.
          \item It is an observational study, because the workers were weighed every week.
          \item It is a census, because every meal was recorded.\hfill\gradedright{\textit{Letter:}~\gradedblank[id=t1, type=choice, points=1, answer=B]{0.45in}{0.2in}}\rule[-0.2in]{0pt}{0.48in}
        \end{enumerate}

  \boxguard[11]
  \item A gym chain randomly selected 300 of the 18{,}000 members at its \textbf{Ohio} locations
        and asked whether they would pay more for longer weekend hours. To which group can the
        result be generalized?
        \begin{enumerate}[label=\textbf{(\Alph*)}, leftmargin=2.2em, itemsep=2pt]
          \item All adults who live in Ohio.
          \item All members of the gym chain, in every state.
          \item Only the 300 members who were asked.
          \item No one, because a sample is not a census.
          \item \ans{All members of the gym chain's Ohio locations.}\hfill\gradedright{\textit{Letter:}~\gradedblank[id=t2, type=choice, points=1, answer=E]{0.45in}{0.2in}}\rule[-0.2in]{0pt}{0.48in}
        \end{enumerate}

  \boxguard[12]
  \item On a holiday weekend an airline randomly selects 6 of its 150 flights and surveys
        \textbf{every} passenger on those 6 flights about the boarding process. Which best
        describes this sample?
        \begin{enumerate}[label=\textbf{(\Alph*)}, leftmargin=2.2em, itemsep=2pt]
          \item A stratified random sample, because the passengers are split into flights.
          \item A simple random sample, because every passenger had a chance to be chosen.
          \item \ans{A cluster sample: everyone on a few random flights is surveyed.}
          \item A stratified random sample, because every passenger on a chosen flight is surveyed.
          \item A convenience sample, because passengers on a plane are easy to reach.\hfill\gradedright{\textit{Letter:}~\gradedblank[id=t3, type=choice, points=1, answer=C]{0.45in}{0.2in}}\rule[-0.2in]{0pt}{0.48in}
        \end{enumerate}

  \boxguard[13]
  \item A teacher's 30 students sit in 6 rows of 5. To choose 5 students, she rolls a die,
        picks the row that matches the number, and takes all 5 students in that row. Every
        student has the same $\tfrac{1}{6}$ chance of being chosen. Is this a simple random
        sample of 5 students from the class?
        \begin{enumerate}[label=\textbf{(\Alph*)}, leftmargin=2.2em, itemsep=2pt]
          \item \ans{No, because two students who sit in different rows can never both be chosen.}
          \item Yes, because every student has the same 1-in-6 chance of being chosen.
          \item Yes, because a die is a fair chance device.
          \item No, because a simple random sample must be chosen with replacement.
          \item No, because the rows are different sizes.\hfill\gradedright{\textit{Letter:}~\gradedblank[id=t4, type=choice, points=1, answer=A]{0.45in}{0.2in}}\rule[-0.2in]{0pt}{0.48in}
        \end{enumerate}

  \boxguard[14]
  \item A school district takes a simple random sample of 300 parents and asks, ``Given the many
        studies showing that later start times improve students' health, do you support moving
        the school day to start at 8:30?'' Which problem is most likely in this survey?
        \begin{enumerate}[label=\textbf{(\Alph*)}, leftmargin=2.2em, itemsep=2pt]
          \item Undercoverage bias: some parents were left off the list.
          \item Voluntary response bias: the parents decide whether to be in the sample.
          \item No bias, because the 300 parents were chosen at random.
          \item \ans{Response bias: the wording leads parents toward ``yes,'' so the estimate is
                probably too high.}
          \item Nonresponse bias, because the parents were chosen at random.\hfill\gradedright{\textit{Letter:}~\gradedblank[id=t5, type=choice, points=1, answer=D]{0.45in}{0.2in}}\rule[-0.2in]{0pt}{0.48in}
        \end{enumerate}

  \boxguard[14]
  \item A radio station asked listeners to call in about a proposed new stadium: 9{,}000 called,
        and 68\% opposed it. A polling firm asked a random sample of 600 of the city's
        registered voters the same question, and 45\% opposed it. Which statement is best?
        \begin{enumerate}[label=\textbf{(\Alph*)}, leftmargin=2.2em, itemsep=2pt]
          \item The call-in poll is more trustworthy, because its sample is 15 times larger.
          \item \ans{The random sample is more trustworthy: the call-in poll has voluntary response
                bias, and a larger sample does not fix bias.}
          \item The two results differ only because of sampling variability.
          \item Averaging the two, about 56\% of voters oppose the stadium.
          \item Neither result is useful, because neither one is a census.\hfill\gradedright{\textit{Letter:}~\gradedblank[id=t6, type=choice, points=1, answer=B]{0.45in}{0.2in}}\rule[-0.2in]{0pt}{0.48in}
        \end{enumerate}

  \boxguard[14]
  \item A school wants to know whether a free after-school tutoring program raises math grades.
        It will randomly assign 90 student volunteers to the tutoring program or to no tutoring
        for one semester. What is the main purpose of the random assignment?
        \begin{enumerate}[label=\textbf{(\Alph*)}, leftmargin=2.2em, itemsep=2pt]
          \item \ans{To make the two groups roughly balanced on variables such as motivation and prior
                grades, so a difference in grades can be traced to the tutoring.}
          \item So the results can be generalized to all students in the state.
          \item To make sure each group has exactly 45 students.
          \item To remove the need for a group that gets no tutoring.
          \item To guarantee that the tutored group earns higher grades.\hfill\gradedright{\textit{Letter:}~\gradedblank[id=t7, type=choice, points=1, answer=A]{0.45in}{0.2in}}\rule[-0.2in]{0pt}{0.48in}
        \end{enumerate}

  \boxguard[15]
  \item To test whether a vitamin~D supplement reduces colds, 200 adults are randomly assigned to
        a vitamin~D pill or to an identical-looking pill with no active ingredient. The adults
        do not know which pill they got, but the nurses who record their colds do. Which
        statement is correct?
        \begin{enumerate}[label=\textbf{(\Alph*)}, leftmargin=2.2em, itemsep=2pt]
          \item The identical pill is a placebo, and the experiment is double-blind.
          \item The identical-pill group is not needed, because the nurses record colds carefully.
          \item The experiment is double-blind because it compares two pills.
          \item The identical pill is a confounding variable.
          \item \ans{It is a placebo, and the study is single-blind.}\hfill\gradedright{\textit{Letter:}~\gradedblank[id=t8, type=choice, points=1, answer=E]{0.45in}{0.2in}}\rule[-0.2in]{0pt}{0.48in}
        \end{enumerate}

  \boxguard[15]
  \item Each of 20 athletes will run a 100-meter sprint twice, on different days: once after a
        sports drink and once after water. Athletes often run faster the second time, once they
        know the course. Which plan is best?
        \begin{enumerate}[label=\textbf{(\Alph*)}, leftmargin=2.2em, itemsep=2pt]
          \item Every athlete drinks water before the first sprint and the sports drink before the
                second.
          \item Let each athlete choose which drink comes first.
          \item \ans{For each athlete, flip a coin: heads, the sports drink comes first; tails, water.}
          \item Each athlete gets both drinks, so no randomization is needed.
          \item Randomly assign 10 athletes to each drink, both days.\hfill\gradedright{\textit{Letter:}~\gradedblank[id=t9, type=choice, points=1, answer=C]{0.45in}{0.2in}}\rule[-0.2in]{0pt}{0.48in}
        \end{enumerate}

  \boxguard[15]
  \item A researcher randomly selects 600 adults from a state's list of registered voters and
        asks how many hours a week they volunteer and how satisfied they are with life. Adults
        who volunteer more report higher satisfaction, and the difference is statistically
        significant. Which conclusion is most appropriate?
        \begin{enumerate}[label=\textbf{(\Alph*)}, leftmargin=2.2em, itemsep=2pt]
          \item Volunteering causes higher life satisfaction for all registered voters in the state.
          \item Volunteering causes higher life satisfaction, but only for these 600 adults.
          \item No conclusion is possible, because this was not an experiment.
          \item \ans{Among the state's registered voters, those who volunteer more tend to be more
                satisfied, but we cannot conclude that volunteering causes it.}
          \item Volunteering causes higher life satisfaction for all adults in the country.\hfill\gradedright{\textit{Letter:}~\gradedblank[id=t10, type=choice, points=1, answer=D]{0.45in}{0.2in}}\rule[-0.2in]{0pt}{0.48in}
        \end{enumerate}

\end{enumerate}

\vspace{0.14in}

% No \newpage here: the guard keeps the Section II strip off a page edge.
\boxguard[30]
\begin{headlinebox}{goldbg}
\small
\textbf{Section II --- Free Response} \hfill \textbf{10 points}\\
Show your reasoning. Answer \emph{in context} --- name the people or things studied and the
variable measured.
\end{headlinebox}

\vspace{0.10in}

\begin{enumerate}[leftmargin=1.9em, itemsep=0.16in, resume]

  \boxguard[26]
  \item \textbf{(4 points)} A school district has 6{,}000 high school students in grades 9
        through 12, in 240 homerooms. It wants to estimate the proportion of its students who
        have a part-time job.

        \vspace{6pt}
        \textbf{(a)} Describe how to use a random number generator to choose a simple random
        sample of 150 students.\\[4pt]
        \gradedhere{\gradedblank[id=t11a, type=rubric, points=1, rubric=../rubrics/apstats-u03-actual-test-t11a.yaml]{\linewidth+6pt}{\gradedlinesheight{3}}}\ansline{Label the students 1 to 6{,}000. Use the generator to produce random}\\
        \ansline{integers from 1 to 6{,}000, skipping any repeat, until 150 different labels}\\
        \ansline{are chosen. Survey the 150 students with those labels.}\\
        \ansline{}\\
        \ansline{}\\
        \ansline{}\\
        \ansline{}\\
        \ansline{}

        \boxguard[16]
        \vspace{6pt}
        \textbf{(b)} \emph{Plan A:} randomly choose 5 of the 240 homerooms and survey every
        student in them. \emph{Plan B:} take a simple random sample of students from each of the
        four grades. Name the sampling method each plan uses, and explain why Plan B should give
        a more precise estimate.\\[4pt]
        \gradedhere{\gradedblank[id=t11b, type=rubric, points=1, rubric=../rubrics/apstats-u03-actual-test-t11b.yaml]{\linewidth+6pt}{\gradedlinesheight{3}}}\ansline{Plan A is a cluster sample; Plan B is a stratified random sample.}\\
        \ansline{Having a job is likely related to grade (older students work more), so each}\\
        \ansline{grade is alike inside. An SRS from every grade represents each one, so the}\\
        \ansline{estimate varies less. Five random homerooms could be mostly one grade.}\\
        \ansline{}\\
        \ansline{}\\
        \ansline{}\\
        \ansline{}

        \boxguard[16]
        \vspace{6pt}
        \textbf{(c)} The district phones the homes of the 150 students in its simple random
        sample between 5 and 8 p.m., and reaches only 96 of them. Name the bias this creates.
        Explain who is likely to be under-represented, and whether the estimate is likely too
        high or too low.\\[4pt]
        \gradedhere{\gradedblank[id=t11c, type=rubric, points=2, rubric=../rubrics/apstats-u03-actual-test-t11c.yaml]{\linewidth+6pt}{\gradedlinesheight{3}}}\ansline{Nonresponse bias: 54 of the 150 chosen students' homes were not reached.}\\
        \ansline{Students who work evening shifts are likely not home from 5 to 8 p.m., so}\\
        \ansline{students with jobs are under-represented. The estimate of the proportion}\\
        \ansline{with a part-time job is likely too low.}\\
        \ansline{}\\
        \ansline{}\\
        \ansline{}\\
        \ansline{}

  \boxguard[24]
  \item \textbf{(6 points)} A grower wants to know whether a new fertilizer increases the yield
        of tomato plants. Thirty-six young plants in her greenhouse will be randomly assigned,
        18 to the new fertilizer and 18 to the standard fertilizer. After ten weeks she weighs
        the tomatoes each plant has produced (kg).

        \vspace{6pt}
        \textbf{(a)} Identify the experimental units, the treatments, and the response
        variable.\\[4pt]
        \gradedhere{\gradedblank[id=t12a, type=rubric, points=1, rubric=../rubrics/apstats-u03-actual-test-t12a.yaml]{\linewidth+6pt}{\gradedlinesheight{3}}}\ansline{Units: the 36 tomato plants.}\\
        \ansline{Treatments: the new fertilizer; the standard fertilizer.}\\
        \ansline{Response: the weight of tomatoes each plant produces in ten weeks (kg).}\\
        \ansline{}\\
        \ansline{}\\
        \ansline{}\\
        \ansline{}\\
        \ansline{}

        \boxguard[14]
        \vspace{6pt}
        \textbf{(b)} Describe how the grower could use 36 identical slips of paper to carry out
        the random assignment.\\[4pt]
        \gradedhere{\gradedblank[id=t12b, type=rubric, points=1, rubric=../rubrics/apstats-u03-actual-test-t12b.yaml]{\linewidth+6pt}{\gradedlinesheight{3}}}\ansline{Number the plants 1 to 36 and write each number on one of 36 identical}\\
        \ansline{slips. Mix them in a hat and draw 18 without replacement: those plants get}\\
        \ansline{the new fertilizer, and the other 18 get the standard fertilizer.}\\
        \ansline{}\\
        \ansline{}\\
        \ansline{}\\
        \ansline{}\\
        \ansline{}

        \boxguard[15]
        \vspace{6pt}
        \textbf{(c)} Half the plants sit on the sunny side of the greenhouse and half on the shady
        side, and sunlight affects yield. Explain how to block on location, and what chance must
        still decide inside each block.\\[4pt]
        \gradedhere{\gradedblank[id=t12c, type=rubric, points=1, rubric=../rubrics/apstats-u03-actual-test-t12c.yaml]{\linewidth+6pt}{\gradedlinesheight{3}}}\ansline{Make two blocks: the 18 plants on the sunny side and the 18 on the shady side.}\\
        \ansline{Inside each block, chance (slips in a hat) assigns 9 plants to the new}\\
        \ansline{fertilizer and 9 to the standard. Compare the fertilizers within each block.}\\
        \ansline{}\\
        \ansline{}\\
        \ansline{}\\
        \ansline{}\\
        \ansline{}

        \boxguard[30]
        \vspace{6pt}
        \textbf{(d)} The new-fertilizer plants produced a mean of \textbf{0.5 kg more} tomatoes
        than the standard plants. A computer reshuffled the 36 yields into two new groups of 18,
        100 times, as if the fertilizer made no difference. How many reshuffles gave a
        difference of 0.5 kg or more? Is the observed difference statistically significant?
        Explain.
        \begin{center}
        \begin{tikzpicture}[x=0.95cm, y=0.1cm]
          \draw[thick] (-5.5,0) -- (5.5,0);
          \foreach \x/\l in {-5/-1.25,-4/-1.00,-3/-0.75,-2/-0.50,-1/-0.25,0/0,
                             1/0.25,2/0.50,3/0.75,4/1.00,5/1.25} {
            \draw (\x,0) -- (\x,-1.4); \node[below, font=\tiny] at (\x,-1.2) {$\l$};}
          \foreach \x/\n in {-5/1,-4/1,-3/5,-2/11,-1/20,0/24,1/20,2/11,3/5,4/1,5/1}
            \foreach \k in {1,...,\n} \fill[navy] (\x,\k) circle (1.3pt);
          \draw[redacc, dashed, thick] (2,12) -- (2,22);
          \node[font=\tiny\bfseries, redacc, above] at (2,22) {observed 0.5};
          \node[font=\tiny] at (0,-7) {Simulated difference in mean yield (kg), new $-$ standard};
          \node[font=\scriptsize\bfseries] at (0,29) {100 random reshuffles};
        \end{tikzpicture}
        \end{center}
        \gradedhere{\gradedblank[id=t12d, type=rubric, points=2, rubric=../rubrics/apstats-u03-actual-test-t12d.yaml]{\linewidth+6pt}{\gradedlinesheight{3}}}\ansline{18 of the 100 reshuffles gave a difference of 0.5 kg or more.}\\
        \ansline{No, it is not significant: a difference this large happens about 18\% of}\\
        \ansline{the time from the luck of the random assignment alone, as if the fertilizer}\\
        \ansline{did nothing. Chance easily explains a 0.5 kg difference.}\\
        \ansline{}\\
        \ansline{}\\
        \ansline{}\\
        \ansline{}

        \boxguard[15]
        \vspace{6pt}
        \textbf{(e)} The grower says, ``The plants were randomly assigned, so this proves the new
        fertilizer increases yield for all tomato plants.'' Explain what is wrong with this
        claim, both about cause and about \emph{all tomato plants}.\\[4pt]
        \gradedhere{\gradedblank[id=t12e, type=rubric, points=1, rubric=../rubrics/apstats-u03-actual-test-t12e.yaml]{\linewidth+6pt}{\gradedlinesheight{3}}}\ansline{Cause: random assignment allows a cause conclusion only for a significant}\\
        \ansline{difference, and 0.5 kg is not, so there is no convincing evidence of any}\\
        \ansline{effect. All plants: these 36 were not randomly selected from all tomato}\\
        \ansline{plants, so a result would apply only to plants like these, in her greenhouse.}\\
        \ansline{}\\
        \ansline{}\\
        \ansline{}\\
        \ansline{}

\end{enumerate}

\end{document}
